\documentclass[10pt,twocolumn]{article}
\usepackage[T1]{fontenc}
\usepackage[utf8]{inputenc}
\usepackage{lmodern}
\usepackage[margin=1.8cm,top=2.4cm,bottom=2cm,columnsep=0.6cm]{geometry}
\usepackage{fancyhdr}
\usepackage{titlesec}
\usepackage{booktabs}
\usepackage{amsmath,amssymb,amsfonts}
\usepackage{microtype}
\usepackage{xcolor}
\usepackage{mdframed}
\usepackage{parskip}
\usepackage{enumitem}
\usepackage{hyperref}

\definecolor{rulered}{RGB}{180,30,30}
\definecolor{boxgray}{RGB}{245,245,245}
\definecolor{keywordgray}{RGB}{100,100,100}

\pagestyle{fancy}
\fancyhf{}
\fancyhead[L]{\textsc{\small Claude Mag}}
\fancyhead[C]{\small\textit{Machine Learning Research Digest}}
\fancyhead[R]{\small 2026-10-28}
\fancyfoot[C]{\small\thepage}
\renewcommand{\headrulewidth}{0.8pt}

\titleformat{\section}{\normalsize\bfseries\color{rulered}}{}{0pt}{}%
  [\vspace{-4pt}\color{rulered}\rule{\columnwidth}{0.4pt}]
\titlespacing{\section}{0pt}{8pt}{4pt}

\hypersetup{colorlinks=true,urlcolor=rulered,linkcolor=black}

\begin{document}
\setlength{\parindent}{0pt}
\setlength{\parskip}{4pt}

{\small\color{keywordgray}\textbf{Keywords:}
  pixel diffusion transformers \textbullet{}
  prediction parameterization \textbullet{}
  MFlow adaptive interpolation \textbullet{}
  ConvRefiner}\par\vspace{4pt}

{\LARGE\bfseries\raggedright
  Rethinking the Design Space of Vanilla Pixel
  Diffusion Transformers}\par\vspace{4pt}

{\large\itshape\raggedright
  MFlow Learns a Per-Direction Interpolation Between $x_0$- and
  Velocity Prediction, Closing the Quality Gap Between Pixel
  and Latent Diffusion}\par\vspace{6pt}

{\small\textbf{By} Soumik Mukhopadhyay, Long Mai,
  Abhinav Shrivastava, Aniruddha Mahapatra}\par\vspace{2pt}

{\small\color{keywordgray}arXiv:2610.23854 \textbullet{}
  Under review, 2026}
\par\vspace{2pt}\noindent{\color{rulered}\rule{\columnwidth}{0.6pt}}\vspace{6pt}

Pixel-space diffusion transformers avoid VAE tokenization artifacts
but have lagged behind latent diffusion in generation quality, partly
because the optimal prediction parameterization---clean image ($x_0$)
versus velocity ($u$)---depends non-trivially on model capacity and
data statistics. MFlow addresses this by introducing per-principal-
direction gains that interpolate between $x_0$- and $u$-prediction,
learned end-to-end with diffusion training. Experiments on ImageNet,
UCF-101 video, and large-scale text-to-image generation confirm that
MFlow consistently outperforms fixed-parameterization baselines and
that larger models systematically prefer velocity prediction. A
lightweight ConvRefiner trained post-hoc with an adversarial loss
further closes the persistent high-frequency gap.

\begin{mdframed}[backgroundcolor=boxgray,linecolor=rulered,linewidth=1pt,
    innerleftmargin=6pt,innerrightmargin=6pt,
    innertopmargin=6pt,innerbottommargin=6pt]
\textbf{\small AT A GLANCE}\par\vspace{4pt}
\begin{description}[leftmargin=0pt,labelsep=4pt,itemsep=2pt,parsep=0pt,
    font=\small\bfseries]
  \item[Problem:] The best prediction parameterization for pixel DiTs
    varies with model capacity and data; the fixed-choice approach
    leaves quality on the table.
  \item[Why it matters:] Pixel DiTs promise simpler, higher-fidelity
    generation than latent diffusion, but their quality gap has
    blocked adoption.
  \item[Method:] Per-direction gain $m_i \in [0,1]$ interpolates
    $x_0$- and $u$-prediction; gains are trained jointly with the
    diffusion model. ConvRefiner handles residual high-frequency
    deficits post-hoc.
  \item[Result:] Competitive or improved FID/FVD vs.\ fixed-
    parameterization baselines on ImageNet, UCF-101, and text-to-image.
  \item[Evidence:] Ablations across model scales, architectures, and
    modalities; systematic $m_i$ analysis validating the capacity-
    preference relationship.
\end{description}
\end{mdframed}

\section{The Big Picture}

Latent diffusion models (LDMs) such as SD3 and FLUX compress images
into a low-dimensional latent code before diffusion, enabling fast
generation but introducing VAE reconstruction artifacts. Pixel
diffusion transformers (pixel DiTs) sidestep this by diffusing
directly in pixel space, preserving sharp boundaries and avoiding
compression bottlenecks.

The quality gap between pixel DiTs and LDMs has been attributed to
less stable training and slower convergence. This paper identifies
an underappreciated cause: the prediction parameterization. The
choice between predicting the clean image ($x_0$), the noise
($\varepsilon$), or the velocity ($u = \alpha x_0 - \sigma\varepsilon$)
determines the gradient signal across noise levels and model capacity
in ways that have not been studied systematically for pixel DiTs.

\section{Why Existing Approaches Fall Short}

\textbf{Fixed $x_0$-prediction (JiT):}
JiT showed that $x_0$-prediction outperforms $\varepsilon$-prediction
for pixel DiTs at moderate scale, but this advantage diminishes or
reverses at larger model sizes. The fixed choice becomes suboptimal
as models scale.

\textbf{Fixed $u$-prediction (flow matching):}
Velocity (flow matching) prediction works well for latent diffusion
but underperforms $x_0$-prediction at smaller pixel DiT scales,
because low-capacity models lack the expressiveness to accurately
predict the velocity at all noise levels.

\textbf{AsymFlow:}
A binary rank-asymmetric assignment of $x_0$ vs.\ $u$ prediction
to different principal components is a special case of MFlow with
step-function gains, and loses the benefit of smooth interpolation.

\section{Core Mechanism}

MFlow parameterizes the prediction target for the $i$-th principal
direction of the data as:
\[
\hat{y}_i = m_i \cdot u_i + (1 - m_i) \cdot (x_0)_i
\]
where the scalar gain $m_i \in [0, 1]$ is the sole additional
parameter per direction. Setting $m_i = 0$ recovers clean-image
prediction; $m_i = 1$ recovers velocity prediction. Gains are
initialized from empirical signal-to-noise statistics and updated
with gradient descent as part of the standard diffusion training loop.

\textbf{ConvRefiner} is a lightweight CNN trained to predict the
residual between a generated pixel DiT image and the real image in
the frequency domain. It is trained post-hoc with a
spectral adversarial loss while keeping the diffusion backbone
frozen, adding only modest compute.

\section{Technical Formulation}

Let $x_0 \in \mathbb{R}^d$ be a clean image in pixel space.
The flow matching diffusion process defines a noisy image:
\[
x_t = \alpha_t x_0 + \sigma_t \varepsilon, \quad \varepsilon \sim \mathcal{N}(0, I)
\]
The velocity target is $u = \alpha_t x_0 - \sigma_t \varepsilon$.
Under MFlow, the model $f_\theta$ predicts:
\[
\hat{y}_i(x_t, t) = m_i f_\theta^{(u)}(x_t, t)_i + (1 - m_i) f_\theta^{(x_0)}(x_t, t)_i
\]
The training loss is the standard flow-matching objective:
\[
\mathcal{L}_\mathrm{MFlow} = \mathbb{E}_{t, x_0, \varepsilon}
  \left[\sum_i \left\|\hat{y}_i(x_t, t) - y_i^\mathrm{target}\right\|^2\right]
\]
where $y_i^\mathrm{target} = m_i u_i + (1 - m_i)(x_0)_i$.
Gradients flow through both $f_\theta$ and $\{m_i\}$.

The ConvRefiner adds:
\[
\hat{x}_\mathrm{refined} = G_\mathrm{ref}(\hat{x}_0) + \hat{x}_0
\]
with $G_\mathrm{ref}$ trained to minimize a spectral GAN loss between
$\hat{x}_\mathrm{refined}$ and real images.

\section{Learning or Inference Procedure}

\textbf{Training:} Standard diffusion training with an extra gradient
term for $\{m_i\}$ at each step. The gains are initialised as
$m_i = \text{SNR}_i / (1 + \text{SNR}_i)$ where $\text{SNR}_i$ is
the signal-to-noise ratio along principal component $i$ at the median
noise level. ConvRefiner is trained in a second stage with the
backbone frozen.

\textbf{Inference:} Identical to standard flow-matching sampling
(DDIM, DPM-Solver, or Euler method). The gain $m_i$ does not affect
inference speed. ConvRefiner adds one forward pass with the lightweight
CNN after the diffusion chain.

\section{What the Guarantee Says}

The MFlow loss is a strict generalization of fixed $x_0$- and
$u$-prediction: setting all $m_i = 0$ recovers the $x_0$ loss,
setting all $m_i = 1$ recovers the velocity loss. Therefore,
MFlow can only improve on either fixed choice---its optima
weakly dominate---assuming sufficient optimization.

For capacity scaling, the paper proves that for a linear model,
the optimal $m_i^*$ satisfies:
\[
m_i^* = \frac{\mathrm{Var}[u_i]}{\mathrm{Var}[u_i] + \mathrm{Var}[(x_0)_i]}
\]
which is an increasing function of the intrinsic dimensionality of
the $i$-th component---larger models, which can represent more
directions, benefit more from velocity prediction.

\section{Experimental Findings}

\begin{itemize}[itemsep=2pt,parsep=0pt]
  \item \textbf{ImageNet 256$\times$256:} MFlow reduces FID by
    $15$--$20\%$ over fixed $x_0$- and $u$-prediction baselines
    at the same model scale.
  \item \textbf{UCF-101 video:} Consistent FVD improvement across
    all model sizes; gains are largest for the largest model ($+18\%$
    over velocity-only baseline).
  \item \textbf{Large-scale text-to-image:} MFlow at 1B parameters
    matches or exceeds the prior best pixel DiT at 3B parameters.
  \item \textbf{ConvRefiner:} Adds 1.2\,ms per image and reduces
    high-frequency FID deficit by 40\% without changing global FID.
  \item \textbf{Capacity scaling:} Empirical $m_i$ values shift from
    near-zero to near-one as model size increases, confirming the
    theoretical prediction.
\end{itemize}

\section{Ablations and Interpretation}

\textbf{Gain initialization:}
Random initialization of $m_i$ converges to similar final values
but requires $2\times$ more training steps; SNR-based initialization
is strongly recommended.

\textbf{Per-direction vs.\ global gain:}
A single shared gain $m$ for all directions reduces FID by only 60\%
of the per-direction improvement, confirming that different principal
components benefit differently.

\textbf{ConvRefiner training signal:}
Replacing the spectral GAN loss with an $\ell_2$ loss in pixel space
fails to recover high-frequency detail---the adversarial objective
is essential.

\section*{Reference}

Soumik Mukhopadhyay, Long Mai, Abhinav Shrivastava, Aniruddha Mahapatra.
\textbf{Rethinking the Design Space of Vanilla Pixel Diffusion Transformers}.
arXiv:2610.23854. Under review, 2026.
\url{https://arxiv.org/abs/2610.23854}

\end{document}
